\addcontentsline{toc}{section}{Abstract}
\vspace*{3cm}
\noindent \large \textbf{\Huge Abstract}
\\[1.5cm]
Passive Radar, also referred to as Bistatic Radar, Passive Coherent Location, or Non-Cooperative Radar, utilizes existing transmitter infrastructure to detect targets based on their skin returns. The transmitters, also known as illuminators of opportunity, are uncooperative transmitters that are not explicitly dedicated for Passive Radar. Examples of such transmitters include Digital Audio Broadcasting, Frequency Modulated radio, Digital Video Broadcasting Terrestrial/Second Generation, Digital Video Broadcasting Satellite/Second Generation and cellular base transmitters.
\newline
\newline
Using FM Passive Radar to detect and visually follow the trajectories of targets presents several challenges, such as clutter, direct and multi-path interference, and poor range resolution. However, it also offers advantages, such as long integration times that allow for high Doppler resolution. Although extensive research has been conducted on tracking filters and their applications across multiple domains, there has been limited exploration on the application of various tracking filters to track commercial aircraft using FM radio signals.
\newline
\newline
 This research aims to explore, develop and implement various tracking filters that can track commercial aircraft using  Passive Radar data. The tracking filters' performance will be evaluated based on several parameters, including computational load, data association and the log-likelihood as a statistical benchmark. By investigating and implementing multiple tracking filters for Passive Radar, this project will contribute towards the advancement of Passive Radar tracking systems.